% P17 discussion: evidence-bound, no new experiment, no causal mechanism claim
\chapter{Bàn luận về thiết kế thực nghiệm và tính hợp lệ}
\label{chap:ban-luan}

\section{Tại sao điểm số tổng thể chưa đủ}
Trong NIDS, việc chọn một pipeline chỉ theo Macro-F1 hoặc Accuracy có thể dẫn tới hiểu nhầm mục tiêu sử dụng. Official TEST chứa 37\,000 dòng Normal và 45\,332 dòng Attack. Với cùng tập này, Stacking-42 giảm số FP từ 10\,003 xuống 9\,861 nhưng tăng FN từ 734 lên 852 so với RF-42. Nếu người vận hành đặt chi phí FN lớn hơn nhiều chi phí FP thì đánh đổi này có thể không mong muốn; trong một môi trường quá tải bởi cảnh báo giả, việc giảm FP lại có giá trị riêng. Tuy nhiên, nghiên cứu không đo chi phí thực tế nên chỉ mô tả đánh đổi, không gọi một pipeline là tốt nhất cho mọi tình huống.

Số ít trường hợp Stacking sửa lỗi của RF và số trường hợp ngược lại còn cho phép nhìn sâu hơn vào kết luận thống kê. Trong bảng paired correctness có 71\,104 mẫu hai mô hình đều đúng, 10\,222 mẫu đều sai, 491 mẫu RF-42 đúng nhưng Stacking-42 sai và 515 mẫu Stacking-42 đúng nhưng RF-42 sai. Sự cân bằng tương đối giữa hai ô sai khác này giải thích vì sao kiểm định về \emph{tính đúng tổng thể} không cho thấy một ưu thế rõ, dù FP/FN theo từng lớp đổi theo những hướng khác nhau.

\begin{table}[htbp]\centering
\caption{Tính đúng/sai của RF-42 và Stacking-42 theo từng cặp dự đoán, official TEST.}
\label{tab:paired-correctness}
\begin{tabular}{p{0.36\linewidth}rr}\toprule
 & Stacking đúng & Stacking sai\\\midrule
RF đúng & 71\,104 & 491\\
RF sai & 515 & 10\,222\\\bottomrule
\end{tabular}
\end{table}

Các ô tính đúng/sai chỉ chứa nhãn đúng/sai tổng thể, trong khi confusion matrix còn phân biệt Normal/Attack. Vì vậy McNemar về hai ô discordant không thể thay thế kiểm tra FPR/FNR hoặc Macro-F1. Việc một phép kiểm định không bác bỏ giả thuyết không đồng nghĩa chứng minh hai mô hình tương đương; nó chỉ phản ánh sức mạnh bằng chứng trong phép thử và dữ liệu được dùng.

\section{Sai số theo nhóm tấn công}
Trên 6\,062 mẫu Fuzzers thật, RF-42 bỏ sót 621 trong khi Stacking-42 bỏ sót 754, tức Stacking tăng 133 FN. Trên Generic có 18\,871 mẫu, RF-42 bỏ sót 57 và Stacking-42 bỏ sót 14, nghĩa là Stacking giảm 43 FN. Hai thay đổi ngược hướng cho thấy một metric gộp toàn bộ Attack có thể che ảnh hưởng không đồng đều giữa các nhóm.

Vì mô hình được huấn luyện cho binary Normal/Attack, đây là phân nhóm những nhãn thật phục vụ phân tích hậu nghiệm. Không được diễn giải các đếm FN thành bảng multiclass accuracy. Đặc trưng \texttt{sttl}, \texttt{proto}, \texttt{state} và các biến khác có thể đồng xuất hiện với mẫu bị dự đoán khác nhau, nhưng việc cùng xuất hiện chưa xác định feature là nguyên nhân. Muốn kiểm tra cơ chế cần một đối chứng có thể can thiệp và không làm thay đổi định nghĩa mẫu hoặc vô tình nhìn lại official TEST trong khâu thiết kế.

\section{Vai trò và giới hạn của Feature Selection}
Hai track 42 và 17 feature được so sánh dưới những cấu hình final khác nhau. RF-42 dùng 300 trees và \texttt{min\_samples\_leaf}=1; RF-17 dùng 150 trees và \texttt{min\_samples\_leaf}=2. Khi RF-17 có thời gian train thấp hơn, lợi ích quan sát được có thể do cả số cây, kích thước lá, số feature lẫn các tương tác giữa chúng. Một thí nghiệm ablation chuẩn để tách nguyên nhân phải giữ các tham số còn lại cố định, dùng cùng seed/protocol và lặp lại việc chọn feature bên trong vòng đánh giá.

Trong dự án, stable 17 features được chọn từ toàn TRAIN trước development CV. Official TEST không được dùng trực tiếp để ra danh sách, nhưng điểm CV sau bước chọn toàn TRAIN có thể vẫn lạc quan vì phép chọn đã thấy nhãn của một số fold trước khi chúng làm validation. Đây là giới hạn nghiêm túc, cần đặt ngay cạnh các kết luận về feature selection, không chỉ chép ở phần cuối. Dữ liệu đã lưu không đủ chứng minh rằng một quy trình nested được triển khai đầy đủ.

\section{Sự khác biệt giữa kiểm định, bootstrap và độ ổn định}
Paired bootstrap tạo phân bố hiệu số từ cùng các dòng TEST, giữ quan hệ ghép cặp của hai model. Đối với Macro-F1, CI 95\% của Stacking-42 trừ RF-42 là [-0,000312; 0,001309]; với MCC là [-0,001938; 0,001140]. Hai khoảng chứa 0. Trong khi đó, CI FPR nằm hoàn toàn phía âm và CI FNR nằm hoàn toàn phía dương. Bốn nhận xét phải được giữ cùng nhau: bằng chứng ủng hộ sự thay đổi loại lỗi trên benchmark, không khẳng định một chiến thắng về hiệu năng tổng quát.

Một phép bootstrap trên dòng không tái tạo sự ngẫu nhiên của quá trình tối ưu mô hình. Để đo ổn định do train, cần nhiều lần fit độc lập, có quản lý seed và khả năng so sánh đúng điều kiện compute. Để đo tổng quát hóa giữa các mạng, cần nhiều benchmark độc lập hoặc thiết kế theo thời gian/session. Không được gọi CI hiện tại là độ tin cậy triển khai ngoài dữ liệu hoặc dùng nó thay cho quá trình tái huấn luyện mô hình.

\section{Giới hạn tái lập đầy đủ}
P15 đã có bằng chứng về checksum dữ liệu local, kiểm tra protocol, khôi phục aggregate postprocessing từ dự đoán đã khóa và smoke fit/serialize/reload trên một phần TRAIN. Những phép kiểm tra này bảo vệ phần mềm và quá trình phân tích, nhưng môi trường Python và thư viện tại lần huấn luyện lịch sử chưa được lưu trực tiếp. Bản model cuối ở dạng serialized cũng không được xác minh tồn tại. Như vậy chưa thể tạo một inference demo đảm bảo đúng pipeline cuối từ những file smoke hiện có.

Các model đánh dấu \texttt{SMOKE\_ONLY} được tạo từ tập con nhỏ nhằm kiểm tra khả năng xây dựng, lưu và tải mô hình. Sử dụng chúng làm demo nhận diện cuối rồi gắn kết quả official TEST sẽ làm sai quan hệ giữa model artifact và số liệu. Trong báo cáo, cần giữ phân biệt ba tầng: mã nguồn có thể chạy, phép phân tích aggregate có thể replay và huấn luyện mô hình cuối có thể tái dựng. Hai tầng đầu đã có bằng chứng theo phạm vi, tầng cuối vẫn là gate mở.

\section{Tính hợp lệ bên ngoài và định hướng nghiên cứu tiếp}
Official UNSW-NB15 TEST không phải một benchmark external độc lập với quá trình hình thành bộ dữ liệu. Novel-vector view là tập con của chính official TEST, không đại diện phép kiểm tra temporal split, day split hay host-separated validation. Vì vậy việc tất cả model mất một phần Macro-F1 trên novel-vector view chỉ chứng minh độ nhạy theo định nghĩa chọn dòng đó, không mô tả lỗi trên traffic tương lai.

Một hướng tiếp tục đáng tin cậy là khóa trước các câu hỏi mới, định nghĩa đơn vị độc lập (session/host/time), huấn luyện trên một split có thể kiểm chứng, đánh giá external nếu có quyền truy cập và đối chiếu các số đo vận hành như FPR theo khối lượng traffic. Những thí nghiệm này phải được đặt thành nghiên cứu kế tiếp có train/test mới, không sửa frozen predictions hay threshold của nghiên cứu hiện tại. Việc công khai các giới hạn giữ cho báo cáo có giá trị học thuật ngay cả khi không có cải thiện rõ rệt của Stacking so với RF.

\section{Ánh xạ câu hỏi nghiên cứu với bằng chứng}
\begin{table}[htbp]\centering
\small
\caption{Mức độ trả lời các câu hỏi nghiên cứu bằng artifact hiện có.}
\label{tab:rq-evidence}
\begin{tabularx}{\linewidth}{p{0.11\linewidth}X p{0.26\linewidth}}\toprule
Câu hỏi & Bằng chứng thực có & Ranh giới kết luận\\\midrule
RQ1 & Official paired scores, CI, McNemar & Không xác nhận stable ensemble gain qua repeated training\\
RQ2 & RF-42/RF-17 cost, stable subset & Pipeline comparison, không isolated feature causality\\
RQ3 & Diversity, OOF, voting, stacking & Mô tả liên hệ, không chứng minh cơ chế nhân quả\\
RQ4 & Novel-vector sensitivity & Chưa có temporal/external validation\\
RQ5 & Phân tích imbalance, weight & Chưa có matched controlled resampling\\\bottomrule
\end{tabularx}
\end{table}
